\section{Scope of the formal verification}\label{app:verification}
The Lean development verifies mathematical components under explicit hypotheses. It is not a construction of the entire stochastic foundation from first principles, nor a verification of the figures. The source archive distinguishes finite algebra, conditional stochastic identities, and paper arguments connecting them.

For the jump results, conditional transform identities, Gaussian characterization, the two-point bound, and the tilted-law example are proved. The general conditional-law argument takes a regular conditional distribution with its defining property as input. Numerical evaluation of the small profile gap is outside Lean.

For the realization result, finite Hankel rank, the Gaussian rank obstruction, recurrence arguments, and state identities are proved. Formal necessity requires only a state at a single time and locally Lipschitz evaluations at finitely many maturities. This weaker hypothesis suffices for the dynamic obstruction. Stochastic identities use supplied calculus interfaces; the sufficient stochastic-scale construction assumes a supplied adapted scale process with the required regularity. The general-shape criterion assembles several components. No instance constructing the entire calculus on a nontrivial Brownian basis is supplied by these arguments.

For approximation, the deterministic coefficient bounds, mean-square bounds under the calculus inputs, Gaussian moment and price estimates, measurable maturity versions for the split quasi-exponential kernel, and identification with the recurrent model are proved. The probabilistic estimates assume the specified driver is Brownian. The numerical recurrence choice and numerical examples are not formalized. The revised article also sharpens the original price constants and gives pointwise error bounds by the explicit elementary proof in Section 5; these sharpened constants have not been added to Lean. The formal assembly retains the original, looser constants.

For the splice, the effective-level identity, aggregate step restrictions, diagonalizable case, time--probability quantifiers, and the converse back to the integrated drift equation with explicit drifts are proved. For the fixed-matrix criterion, the block is linear in its coordinates, so its drift and volatility are written explicitly. The core theorem takes the integrated drift condition as an explicit hypothesis rather than adding a stochastic-calculus axiom. The converse is an equation for coefficients, not a construction of the block process.

The local-integrability extension proves pathwise integrability of the explicit converse drifts and maturity-uniform envelopes. The construction of front-end processes as stochastic integrals and progressive measurability of the assembled representation remain paper arguments. No true-martingale theorem follows from these local bounds alone.

For the three-exponent variance example (original Proposition 27, Claim 025), the lab derives the constructed scalar state, stopped-system uniqueness, variance identities, independence needed for the nonzero short-rate jump, and true discounted-bond martingales. The formal result is conditional on compatible natural- and joint-filtration It\^o calculus structures with the same Brownian driver, and their audited SDE, predictability, integral-approximation, and exponential-martingale inputs. Their existence on a nontrivial Brownian basis is not constructed. Under those inputs the model conclusions are proved, rather than assumed. The article's standard-calculus argument and the formal conditional scope must be distinguished.

For varying exponents (original Proposition 46, Claim 042), the parameter derivatives, pointwise residual reduction, and almost-everywhere passage are formalized. The inherited exponent restrictions are fields of the audited upstream interface for the cited classical theorem, AX-16. The result requires positive exponents almost everywhere, including where their coefficient polynomials vanish. It does not treat correlated moving exponents or exponent-zero components. The fixed-decay Svensson corollary (original Proposition 40, Claim 038) is formalized as a pointwise covariance and drift classification; no process existence is asserted. The new offline coefficient checks do not extend the Lean coverage.

The stochastic interfaces contain cited results as hypotheses. Some supporting interfaces elsewhere in the development have only degenerate or vacuous instances with separate convention checks. These are not evidence of a nontrivial Brownian model satisfying every interface. A conditional Lean theorem must not be described as removing the existence assumptions of the stochastic theorem.

The companion source map records original proposition numbers and principal Lean assembly files. It belongs to the reproducibility material; the article's mathematical numbering is independent of the order of the research.
